% !TEX root = ../main.tex

\begin{abstract}
\noindent
Todo el mundo debería tener la oportunidad de programar igual que lo hace cualquier programador, sin importar la condición física, cognitiva o social.
El lenguaje de programación \sculpt (\textit{Simple Cubic Language for Programming Tasks}) es un lenguaje de programación y ecosistema artistico que pretende retar las convenciones de los lenguajes de programación tradicionales y nociones sobre el pensamiento computacional.
\sculpt es un lenguaje de programación visual, tangible y modular el cual disrumpe la interfaz estandar de teclado y pantalla mediante una interfaz física que se puede tocar, manipular y compartir junto a su especificación técnica, código fuente y documentación.
Esto se realiza mediante la construcción de constructos artísticos los cuales se pueden ensamblar para crear esculturas que generan programas computables.
Junto al lenguaje, se presenta \sculpter para crear y ejecutar programas antes de ensamblarlos en \sculpt. Adicionalmente, cuenta con documentación sobre su uso y desarrollo y tutoriales para aprender a programar con \sculpt.
El lenguaje construido es totalmente capaz y Turing-completo el cual sirve como una herramienta creativa tambien.
Usando estas herramientas, se realiza una validación empirica de la experiencia con novatos y expertos en programación.
Los resultados muestran ligera proeza computacional al resolver problemas de programación entre las poblaciones y unas piezas de arte interesantes.
\end{abstract}

\medskip
\noindent\textbf{Palabras clave:} pensamiento computacional, lenguajes de
programaci\'on visuales, dise\~no inclusivo, computaci\'on creativa.

\endinput
